\section{Programas em assembly}
\label{sec:programas}

Três programas foram escritos para demonstrar e testar a CPU. O primeiro, a
sequência de Fibonacci, é o programa de demonstração carregado por padrão na
RAM. Os outros dois são testes \emph{autoverificáveis}: o próprio programa
confere cada resultado e informa no LCD se passou ou falhou, o que permite
testar a CPU diretamente na placa.

\subsection{Sequência de Fibonacci}

O programa da Listagem~\ref{lst:fib} escreve em $MEM[255]$ --- e portanto na
linha 2 do LCD --- os termos da sequência de Fibonacci que cabem em 8 bits:
0, 1, 1, 2, 3, 5, 8, 13, 21, 34, 55, 89, 144 e 233. O termo seguinte,
$144 + 233 = 377$, não cabe em 8 bits; a soma gera \emph{carry}
(OVERFLOW = 1) e o programa termina.

\lstinputlisting[language=asm8, caption={Programa \texttt{fibonacci.asm}.}, label={lst:fib}]{../programs/fibonacci.asm}

\paragraph{Funcionamento.} \sig{Ra} e \sig{Rb} guardam dois termos
consecutivos, $F(n-1)$ e $F(n)$. A cada volta do laço, $F(n)$ é mostrado, guardado
temporariamente em \sig{Rc}, e a soma $F(n-1) + F(n)$ é calculada em \sig{Rb};
em seguida \sig{Ra} recebe o antigo $F(n)$. O programa explora dois detalhes
da arquitetura:

\begin{itemize}[nosep]
  \item \instr{mov} e \instr{ldi} \textbf{não alteram as flags}. Por isso a
        flag OVERFLOW calculada pelo \instr{add} continua valendo quando o
        \instr{bcs} é executado, três instruções depois.
  \item Nos saltos condicionais, o registrador contém o \textbf{endereço de
        destino}: \instr{ldi Rd, fim} carrega o endereço do rótulo \sig{fim}
        (16), e \instr{bcs Rd} salta para lá se houve \emph{carry}.
\end{itemize}

\paragraph{Código de máquina.} A Tabela~\ref{tab:fibbin} mostra o programa
montado, exatamente como fica no bloco de inicialização da RAM. As posições 1,
3, 5, 7, 12 e 15 são o segundo byte das instruções \instr{ldi}, \instr{st} e
\instr{jmp}, e nunca são executadas.

\begin{table}[H]
\centering
\caption{Fibonacci em código de máquina.}
\label{tab:fibbin}
\small
\begin{tabular}{r l l @{\hspace{1cm}} r l l}
\toprule
End. & Binário & Conteúdo & End. & Binário & Conteúdo \\
\midrule
0 & \texttt{10000011} & \instr{ldi Ra,} & 9  & \texttt{00000100} & \instr{add Rb, Ra} \\
1 & \texttt{00000000} & \quad 0         & 10 & \texttt{10110010} & \instr{mov Ra, Rc} \\
2 & \texttt{10000111} & \instr{ldi Rb,} & 11 & \texttt{10001111} & \instr{ldi Rd,} \\
3 & \texttt{00000001} & \quad 1         & 12 & \texttt{00010000} & \quad 16 (\sig{fim}) \\
4 & \texttt{10000010} & \instr{st Ra,}  & 13 & \texttt{11011100} & \instr{bcs Rd} \\
5 & \texttt{11111111} & \quad 255       & 14 & \texttt{11000000} & \instr{jmp} \\
6 & \texttt{10000110} & \instr{st Rb,}  & 15 & \texttt{00000110} & \quad 6 (\sig{loop}) \\
7 & \texttt{11111111} & \quad 255       & 16 & \texttt{11111111} & \instr{halt} \\
8 & \texttt{10111001} & \instr{mov Rc, Rb} & & & \\
\bottomrule
\end{tabular}
\end{table}

O programa executa 94 instruções até o \instr{halt}, ou seja, cerca de
\SI{188}{\second} na placa. A execução completa, passo a passo, está na
Seção~\ref{sec:passo}.

\subsection{Testes autoverificáveis}

Os programas \sig{test\_alu.asm} e \sig{test\_mem.asm} exercitam todas as
instruções da ISA. Cada teste executa uma operação, confere o resultado e as
flags usando os próprios saltos condicionais, e desvia para \sig{fail} na
primeira divergência. Quando um teste passa, seu número é escrito em
$MEM[255]$; ao final, $MEM[255] = 100$ indica que tudo passou, e
$MEM[255] = 238$ indica falha --- nesse caso, o último número exibido antes do
238 aponta o último teste aprovado.

Para comparar um registrador com um valor esperado, os testes usam a própria
ALU: \instr{ldi Rc, esperado} seguido de \instr{and Rc, Rx} faz a ALU comparar
\sig{Rc} com \sig{Rx} e acender EQUAL se forem iguais, e então
\instr{bneq} desvia para \sig{fail}. A Listagem~\ref{lst:t1} mostra o primeiro
teste.

\begin{lstlisting}[language=asm8, caption={Teste T1 de \texttt{test\_alu.asm}: soma com \emph{carry}.}, label={lst:t1}]
        ldi Ra, 200
        ldi Rb, 100
        add Ra, Rb          ; 200 + 100 = 300 -> 44 com carry
        ldi Rd, fail
        bcc Rd              ; ha carry -> nao pode saltar
        blt Rd              ; 200 < 100 ? nao
        bz  Rd              ; resultado != 0
        bneg Rd             ; 44 e positivo
        ldi Rd, t1a
        bcs Rd              ; deve saltar
        jmp fail
t1a:    ldi Rd, t1b
        bgt Rd              ; 200 > 100 -> salta
        jmp fail
t1b:    ldi Rc, 44
        and Rc, Ra          ; EQUAL = (44 == Ra)
        ldi Rd, fail
        bneq Rd
        ldi Ra, 1
        st  Ra, 0xFF        ; teste 1 aprovado
\end{lstlisting}

Cada condição é testada nos dois sentidos: um salto que \emph{não deve}
ocorrer aponta para \sig{fail}, e um salto que \emph{deve} ocorrer é seguido de
\instr{jmp fail}. Assim, um salto que nunca é tomado e um salto que é sempre
tomado são ambos detectados. A Tabela~\ref{tab:testes} resume o que cada teste
cobre.

\begin{table}[H]
\centering
\caption{Testes autoverificáveis.}
\label{tab:testes}
\small
\begin{tabularx}{\linewidth}{c l X}
\toprule
Teste & Programa & O que verifica \\
\midrule
T1 & \sig{test\_alu} & \instr{add} com \emph{carry} ($200+100=44$): OVERFLOW = 1, GREATER = 1; \instr{bcc}, \instr{blt}, \instr{bz}, \instr{bneg} não tomados; \instr{bcs}, \instr{bgt} tomados \\
T2 & \sig{test\_alu} & \instr{sub} com \emph{borrow} ($3-5=254$): OVERFLOW = 1, NEGATIVE = 1, SMALLER = 1; \instr{bneg}, \instr{blt} tomados \\
T3 & \sig{test\_alu} & resultado zero e igualdade ($77-77$): ZERO = 1, EQUAL = 1, OVERFLOW = 0; \instr{bz}, \instr{beq}, \instr{bcc} tomados \\
T4 & \sig{test\_alu} & \instr{inc} de 255 (0 com \emph{carry}), \instr{dec} de 0 (255 com \emph{borrow}) e casos sem estouro \\
T5 & \sig{test\_alu} & \instr{and}, \instr{or}, \instr{xor}, \instr{not} (valores e OVERFLOW = 0); \instr{not} de 255 gera ZERO \\
T6 & \sig{test\_alu} & \instr{rol}, \instr{ror}, \instr{lsl}, \instr{lsr} sobre \texttt{0x81}, inclusive o bit que sai em OVERFLOW \\
T7 & \sig{test\_mem} & \instr{str}, \instr{st} (escrita direta) e \instr{ldr} (leitura indireta) \\
T8 & \sig{test\_mem} & \instr{push}/\instr{pop}: posição e ordem na pilha, $SP$ iniciado em 254 \\
T9 & \sig{test\_mem} & \instr{jmpr}; \instr{mov}, \instr{ldi} e \instr{st} não alteram flags; sete opcodes inválidos executam como \instr{nop} sem alterar registradores nem flags \\
\bottomrule
\end{tabularx}
\end{table}

Os opcodes inválidos do teste T9 são inseridos no programa como bytes de dados
(\texttt{0x22}, \texttt{0x23}, \texttt{0x5D}, \texttt{0xC4}, \texttt{0xEF},
\texttt{0xF5} e \texttt{0xF0}), cobrindo cada caso de código não usado listado
no enunciado.
